\documentclass[10pt,a4paper]{amsart}
\usepackage{amsmath,amssymb,amsthm,graphicx}
% Retained preamble for the Stacks Project chapter {CH} (shared project preamble).
% The chapter file pinned-original-source.tex only carries \input{{preamble}} on its first line,
% so no preamble text lives inside the chapter itself.  The Stacks Project uses one single shared
% preamble for every chapter: the snapshot stages it three times (sources/stacks-chow/preamble.tex,
% sources/stacks-etale-cohomology/preamble.tex and focused-seeds/AGP/sources/stacks-snapshot-20261011/
% preamble.tex) and those copies are identical in size.  The block below is copied mechanically,
% line for line, from sources/stacks-chow/preamble.tex lines 14-18 and 154-217: the Xy-pic setup and
% the native theorem declarations together with the mathematical macros.  Nothing was added, removed
% or redefined; the wrapper class and the amsmath/amssymb/amsthm/graphicx base packages are supplied
% by the assembler.
% For commutative diagrams we use Xy-pic
% Chapter: stacks-spaces-more-morphisms.tex -- More on Morphisms of Spaces.
\usepackage[all]{xy}

% We use 2cell for 2-commutative diagrams.
\xyoption{2cell}

% Theorem environments.
%
\theoremstyle{plain}
\newtheorem{theorem}[subsection]{Theorem}
\newtheorem{proposition}[subsection]{Proposition}
\newtheorem{lemma}[subsection]{Lemma}

\theoremstyle{definition}
\newtheorem{definition}[subsection]{Definition}
\newtheorem{example}[subsection]{Example}
\newtheorem{exercise}[subsection]{Exercise}
\newtheorem{situation}[subsection]{Situation}

\theoremstyle{remark}
\newtheorem{remark}[subsection]{Remark}
\newtheorem{remarks}[subsection]{Remarks}

\numberwithin{equation}{subsection}

% Macros
%
\def\lim{\mathop{\mathrm{lim}}\nolimits}
\def\colim{\mathop{\mathrm{colim}}\nolimits}
\def\Spec{\mathop{\mathrm{Spec}}}
\def\Hom{\mathop{\mathrm{Hom}}\nolimits}
\def\Ext{\mathop{\mathrm{Ext}}\nolimits}
\def\SheafHom{\mathop{\mathcal{H}\!\mathit{om}}\nolimits}
\def\SheafExt{\mathop{\mathcal{E}\!\mathit{xt}}\nolimits}
\def\Sch{\mathit{Sch}}
\def\Mor{\mathop{\mathrm{Mor}}\nolimits}
\def\Ob{\mathop{\mathrm{Ob}}\nolimits}
\def\Sh{\mathop{\mathit{Sh}}\nolimits}
\def\NL{\mathop{N\!L}\nolimits}
\def\CH{\mathop{\mathrm{CH}}\nolimits}
\def\proetale{{pro\text{-}\acute{e}tale}}
\def\etale{{\acute{e}tale}}
\def\QCoh{\mathit{QCoh}}
\def\Ker{\mathop{\mathrm{Ker}}}
\def\Im{\mathop{\mathrm{Im}}}
\def\Coker{\mathop{\mathrm{Coker}}}
\def\Coim{\mathop{\mathrm{Coim}}}

% Boxtimes
%
\DeclareMathSymbol{\boxtimes}{\mathbin}{AMSa}{"02}

%
% Macros for moduli stacks/spaces
%
\def\QCohstack{\mathcal{QC}\!\mathit{oh}}
\def\Cohstack{\mathcal{C}\!\mathit{oh}}
\def\Spacesstack{\mathcal{S}\!\mathit{paces}}
\def\Quotfunctor{\mathrm{Quot}}
\def\Hilbfunctor{\mathrm{Hilb}}
\def\Curvesstack{\mathcal{C}\!\mathit{urves}}
\def\Polarizedstack{\mathcal{P}\!\mathit{olarized}}
\def\Complexesstack{\mathcal{C}\!\mathit{omplexes}}
% \Pic is the operator that assigns to X its picard group, usage \Pic(X)
% \Picardstack_{X/B} denotes the Picard stack of X over B
% \Picardfunctor_{X/B} denotes the Picard functor of X over B
\def\Pic{\mathop{\mathrm{Pic}}\nolimits}
\def\Picardstack{\mathcal{P}\!\mathit{ic}}
\def\Picardfunctor{\mathrm{Pic}}
\def\Deformationcategory{\mathcal{D}\!\mathit{ef}}
\usepackage{cleveref}
\crefname{originallemma}{original lemma}{original lemmas}
\Crefname{originallemma}{Original lemma}{Original lemmas}
\makeatletter
\newlabel{lemma-exact}{{08B5}{source}{}{sourcepointer-1}{}}
\newlabel{lemma-exact@cref}{{[originallemma][0][]08B5}{[1][0][]source}{}{}{}}
\makeatother
\crefname{originalremark}{original remark}{original remarks}
\Crefname{originalremark}{Original remark}{Original remarks}
\makeatletter
\newlabel{remark-reformulate-existence-theorem}{{08BF}{source}{}{sourcepointer-2}{}}
\newlabel{remark-reformulate-existence-theorem@cref}{{[originalremark][0][]08BF}{[1][0][]source}{}{}{}}
\makeatother
\crefname{originalreference}{original reference}{original references}
\Crefname{originalreference}{Original reference}{Original references}
\makeatletter
\newlabel{spaces-limits-lemma-check-closed-infinitesimally}{{08B2}{source}{}{sourcepointer-3}{}}
\newlabel{spaces-limits-lemma-check-closed-infinitesimally@cref}{{[originalreference][0][]08B2}{[1][0][]source}{}{}{}}
\makeatother
\title{Algebraization of a formal closed subspace of an algebraic space}
\author{The Stacks Project Authors}
\date{stacks-project:spaces-more-morphisms.tex:local-snapshot-20261011; GNU FDL 1.2 or later}
\begin{document}
\maketitle
\noindent\textit{Editorial note.} Counted span is the author's own proof environment, source lines 10168-10191, immediately following the statement at 10145-10166.\par\medskip
\noindent\textit{Editorial note.} Same-name pair note: this is the more general (algebraic spaces) side; the schemes side is not counted beside it.\par\medskip
\noindent\textit{Editorial note.} The Grothendieck existence theorem and its reformulation are same-chapter navigation pointers recorded in assembly.external\_references.\par\medskip
\noindent\textit{Allowed standard prerequisites (editorial).}\begin{itemize}
\item Grothendieck's existence theorem for algebraic spaces. chapter theorem theorem-grothendieck-existence, cited in the proof. Produces the algebraization of the compatible system of ideal sheaves on the thickenings.
\item Reformulation of existence as an equivalence of categories. chapter remark remark-reformulate-existence-theorem, cited in the proof. Gives the form of the theorem applied to the ideals defining the formal closed subspace.
\item Closed subschemes and properness after algebraization. used at the end of the proof (Z proper over S). Reads off that the algebraized object is a closed immersion and proper over S.
\end{itemize}\medskip
\ifdefined\hypertarget\hypertarget{sourcepointer-1}{}\fi
\noindent\textit{Source reference.} Original lemma 08B5: sources/stacks-spaces-more-morphisms/pinned-original-source.tex, source lines 9644--9644. Navigation pointer only: the counted proof cites this statement by \textbackslash{}ref as a standing input; the deduction the counted proof has to make is carried out inside the counted span, and no part of that deduction is delegated to this pointer.\par\medskip
\ifdefined\hypertarget\hypertarget{sourcepointer-2}{}\fi
\noindent\textit{Source reference.} Original remark 08BF: sources/stacks-spaces-more-morphisms/pinned-original-source.tex, source lines 10085--10085. Navigation pointer only: the counted proof cites this statement by \textbackslash{}ref as a standing input; the deduction the counted proof has to make is carried out inside the counted span, and no part of that deduction is delegated to this pointer.\par\medskip
\ifdefined\hypertarget\hypertarget{sourcepointer-3}{}\fi
\noindent\textit{Source reference.} Original reference 08B2: focused-seeds/geometry/S9-spaces-algebraize-formal-closed-subscheme-spaces/ref-tags.txt, source lines 4--4. Navigation pointer only: the counted proof cites this statement by \textbackslash{}ref as a standing input; the deduction the counted proof has to make is carried out inside the counted span, and no part of that deduction is delegated to this pointer.\par\medskip
\begin{lemma}
\label{lemma-algebraize-formal-closed-subscheme}
Let $A$ be a Noetherian ring complete with respect to an ideal $I$.
Write $S = \Spec(A)$ and $S_n = \Spec(A/I^n)$.
Let $X \to S$ be a morphism of algebraic spaces that is separated
and of finite type.
For $n \geq 1$ we set $X_n = X \times_S S_n$.
Suppose given a commutative diagram
$$
\xymatrix{
Z_1 \ar[r] \ar[d] & Z_2 \ar[r] \ar[d] & Z_3 \ar[r] \ar[d] & \ldots \\
X_1 \ar[r]^{i_1} & X_2 \ar[r]^{i_2} & X_3 \ar[r] & \ldots
}
$$
of algebraic spaces with cartesian squares. Assume that
\begin{enumerate}
\item $Z_1 \to X_1$ is a closed immersion, and
\item $Z_1 \to S_1$ is proper.
\end{enumerate}
Then there exists a closed immersion of algebraic spaces $Z \to X$ such that
$Z_n = Z \times_S S_n$ for all $n \geq 1$. Moreover, $Z$ is proper over $S$.
\end{lemma}

\begin{proof}
Let's write $j_n : Z_n \to X_n$ for the vertical morphisms.
As the squares in the statement are cartesian
we see that the base change of $j_n$ to $X_1$ is $j_1$.
Thus Limits of Spaces, Lemma
\ref{spaces-limits-lemma-check-closed-infinitesimally}
shows that $j_n$ is a closed immersion.
Set $\mathcal{F}_n = j_{n, *}\mathcal{O}_{Z_n}$, so that
$j_n^\sharp$ is a surjection $\mathcal{O}_{X_n} \to \mathcal{F}_n$.
Again using that the squares are cartesian we see that
the pullback of $\mathcal{F}_{n + 1}$ to $X_n$ is $\mathcal{F}_n$.
Hence Grothendieck's existence theorem, as reformulated in
Remark \ref{remark-reformulate-existence-theorem},
tells us there exists a map
$\mathcal{O}_X \to \mathcal{F}$
of coherent $\mathcal{O}_X$-modules whose restriction to
$X_n$ recovers $\mathcal{O}_{X_n} \to \mathcal{F}_n$.
Moreover, the support of $\mathcal{F}$ is proper over $S$.
As the completion functor is exact (Lemma \ref{lemma-exact})
we see that $\mathcal{O}_X \to \mathcal{F}$
is surjective. Thus $\mathcal{F} = \mathcal{O}_X/\mathcal{J}$
for some quasi-coherent sheaf of ideals $\mathcal{J}$.
Setting $Z = V(\mathcal{J})$ finishes the proof.
\end{proof}

\end{document}
